\chapter{4}

\section{Montag, 20. August}



Louis griff den Rucksack, erhob sich und schritt entschlossenen in Giorgias Richtung. 
Je weiter er sich vom Feuer entfernte und sich im Dunkel zwischen den
Bäumen vorantastete, desto überzeugter war er. Heute musste es sein.
Seine Augen stellten sich auf die Dunkelheit ein. Schritt für Schritt
schlich er sich zwischen Ästen und Büschen hindurch. Er hob einen
kurzen Stock hoch und schob damit immer wieder einzelne Zweige zur
Seite, bis er den lichteren Teil erreicht hatte. Das Pochen in seinen
Schläfen nahm zu. Es war, als sässe sein Herz direkt im Kopf. Erneut 
diese Angst. Er hoffte, die richtigen Worte zu finden.

Schliesslich sah er Giorgias Silhouette wenige Meter vor sich.
Sie schwankte bedrohlich nahe an der Kante eines langen Steinbrockens.
In kleinen Schritten näherte er sich ihr. Er zögerte. Ihr Gesicht war
nicht klar zu erkennen. Sie schrie. Ein aufs neue Mal. Nun verstand er
die Worte.

\emph{«Vattene da qui. Non voglio più vederti, vaffanculo»}.\footnote{«Hau ab hier. Ich will dich nicht mehr sehen, verpiss dich.»}

Er zuckte erschrocken zusammen. Sie beschimpfte ihn? Warum? Er verstand ihre Schärfe nicht, ihre Wut. Er hatte doch noch gar nichts
gesagt. Wo stand sie da? Tausend Gedanken jagten sich in seinem Kopf.
Und mit einem Mal wallte tiefer Zorn in ihm auf.
Sie hatte kein Recht, ihn zu beschimpfen, nein, das hatte sie
nicht. Ihm, und nur ihm, stand die Wut zu. Louis war, als habe die
Dunkelheit das Zepter übernommen und sich eine boshafte Stimme gewaltsam Zugang zu seinem Gehirn erzwungen. 
Giorgia stand direkt vor dem Abgrund.
Sie gestikulierte wild mit den Armen. Er ging langsam näher auf sie zu. Plötzlich
richtete sie ihre Augen direkt auf ihn. Die Überraschung in ihrem Blick
verstörte ihn. Dann veränderte sich ihr Gesichtsausdruck. Er sah ein
Flehen, wie er es nie zuvor bei ihr gesehen hatte. Er schüttelte den
Kopf. Ihre eben ausgesprochenen Worte passten nicht zu diesem Blick.
Louis presste die Lippen aufeinander. Giorgia streckte ihre Arme, die
Handflächen offen gegen ihn gerichtet, weit von sich. Sie öffnete ihren
Mund.
Louis’ Augen füllten sich mit Tränen der Empörung und ein unaufhaltsamer
Drang, ihre Handgelenke zu fassen, packte ihn.
Jetzt schwankte sie plötzlich nach vorne. Hielt ihm ihre rechte Hand
entgegen. Schlagartig streckte er ihr den Stock hin, sie fasste danach,
er zog, sie zog, beide zogen sie hin und her, dann, ein Schrei, ein
letzter Blick und Giorgia war weg.\\

Vor ihm tat sich nichts als eine gespenstische Leere auf.\\

Mit angehaltenem Atem schlich er sich ganz nach vorne an die Felskante. Sein Körper bebte. Zwischen den dicht hinter ihm
aufgereihten Bäumen drängte sich nun die Musik geradezu auf. Der
gedämpfte Rhythmus kam einer Drohung gleich. Mit geweiteten Augen ging
er in die Knie. Schaute in die Tiefe. Ein dunkles Nichts.
Er starrte auf den Stock in seiner Hand, drückte ihn krampfartig
zusammen und hörte sich flüstern.
«Was war das? Wo ist sie?»
Er kannte den Ort. Er wusste, die Wand war gute dreissig, vielleicht
vierzig Meter hoch. Oft hatten sie darüber Witze gerissen, er und seine
Freunde. «Einmal unten, nie mehr oben» war zum geflügelten Wort
geworden.

Louis liess sich auf den Boden fallen.

Wo ist sie? Wo ist sie?

Er krallte seine Finger in die trockene Erde. Riss die Augen weiter auf.
Nur dunkle Schatten mit dazwischenliegenden helleren Flecken waren zu
erkennen. Er suchte vergeblich nach dem Orange ihres Kleides, wenigstens
als Farbklecks, zwischen Laub und Ästen und glaubte, einzufrieren.\\

Es mochten nur wenige Sekunden vergangen sein. Doch Louis war, als würde er
nach einer ewig langen Absenz wieder zu sich kommen. 
Er erhob sich hastig. Sein Zorn schien wie weggeblasen. 
Er stolperte der Felskante entlang, rannte wie ein getriebenes Tier
seitlich den Hügel hinunter und bog rechts um eine kleine Baumgruppe, 
bis er auf dem schmalen Pfad stand, der parallel der Felswand entlang
führte.

«Giorgia. Giorgia?»
Sein Herz suchte nach dem nächsten Schlag. Aus weiter Ferne war die
Sirene eines Polizeiwagens zu hören. Unvermittelt begannen seine Hände
zu vibrieren. Er hob sie hoch und schaute sie entgeistert an. Es war
chancenlos, die Kontrolle zurückzugewinnen. Er rang nach Luft, folgte kleinschrittig dem Weg und ihn grauste, Giorgia auf dem schmalen Pfad liegend zu finden. Aber nichts war zu sehen. 
Auch der Ort, wo sie hätte aufschlagen müssen, war gähnend leer.

Unmöglich, einfach unmöglich. Wieder und wieder murmelte er die Worte vor sich hin. Sein Kopf zuckte
nach allen Seiten. Nichts, einfach nichts. Keine Giorgia. Er versuchte mehrmals, nach ihr zu rufen. Doch seine Stimme war tot.
Wie sollte er ihr helfen, wenn sie nicht da war? Wo war sie?
Und seine Freunde? Wenn er sie holen würde, was würden sie denken? Es
läge auf der Hand. Es war seine Schuld.

Seine linke Hand hatte sich beruhigt. Die Rechte zitterte weiter und in
ihr der Stock. Dann rannte er panisch davon. Erst später würde er sich damit trösten können, 
dass er nicht mehr hatte denken können. Die Verzweiflung schien sich in seinen Beinen festgebissen zu haben. Sie 
trugen ihn ganz von selbst und rasend schnell davon. 
Warum, warum nur war er Giorgia gefolgt? Er hatte gesehen, wie
zugedröhnt sie war. Er hatte die falsche Gelegenheit genutzt, um noch
einmal - ein einziges Mal - mit ihr zu sprechen. Warum hatte sie ihn
angeschrien, verdammt, was war da gerade passiert?
Er keuchte. Der Schweiss klebte an seinem Körper. Sein Rucksack blieb an
einem Ast hängen. Er strauchelte, fing sich. Dann blieb er mit dem
Rücken an einen Baum gepresst stehen. Stocksteif rang er nach Luft,
schloss die Augen und starrte in sich hinein.

Schliesslich hetzte er ziellos davon.

\sectionbreak

\indentedblock{Damals. Fribourg, sieben.\\
Louis ruft nach Maman, sucht sie in der Wohnung. Er findet sie, schräg
auf dem Sofa hängend, im Wohnzimmer. Stumm, mit geschlossenen Augen
liegt sie da. Sie regt sich nicht. Er schaut erst nur staunend zu ihrem
Gesicht hinunter, sieht dem langen Haar entlang, das den Boden berührt
und stösst schliesslich einen Schrei aus, über den er selbst erschrickt.
«Maman ist tot,» brüllt er, wagt es nicht, sie auch nur zu berühren,
überzeugt, dass sie ihn verlassen hat. Und er rennt aus dem Zimmer in den Flur, raus aus der Wohnung ins 
Treppenhaus, nimmt mehrere Stufen auf einmal, hält sich verkrampft am
Geländer fest, überquert die Strasse vor dem Mehrfamilienhaus, schaut
kaum links noch rechts, stolpert den «Boulevard de Pérolles» entlang und
lässt sich schliesslich leise weinend an der Seitenwand seines «Kiosque»
entlang zu Boden gleiten. Dass er sich dabei den Rücken aufkratzt, spürt
er später. Louis schlottert. Er hört Cathrines Stimme. Sie verabschiedet
sich gerade von einem Kunden. Soll er zu ihr? Seine Lieblingsverkäuferin
würde ihm helfen, stopp, nein, er hört die Worte seines Vaters, nachdem
er ihm einmal freudig erzählt hat, wie lieb sie zu ihm ist, und dass sie
nun seine Schlangenfreundin ist, weil die Süssigkeit so traumhaft
schmeckt und hört wie Vater zu ihm sagt, die sei eine geschwätzige Frau,
und er dürfe ihr nichts Privates sagen, gar nichts - Papa. Er wird zu
Papa gehen. Der wird ihn retten. Rennend erreicht er das nahe gelegene Wohnhaus 
und wischt immer wieder die Augen mit seinem Ärmel trocken.
Louis hofft, dass er da ist. Er kennt das Haus, wo sein Vater immer häufiger wohnt, wenn er nicht bei ihnen ist. Es wird Maman nicht gefallen, dass er ohne zu fragen dahin geht. Sie mag Vaters Kumpel nicht besonders. Aber wie hätte er sie fragen sollen?
Er reckt sich hoch und erinnert sich. Zweiter Knopf von oben rechts, hat Vater gesagt. Er klingelt und wartet. 
Ein Surren erklingt. Louis stösst die Tür auf. Sie ist schwer.}
